


\documentclass[11pt]{article}
\usepackage{amsgen,amsmath,amstext,amsbsy,amsopn,amsfonts,amsthm,amscd}
%\usepackage{bbm}
\usepackage[dvips]{graphicx}
\usepackage{enumerate}
\setlength{\textwidth}{16.0cm} \setlength{\textheight}{25cm}
\voffset=-2.50cm \hoffset=-2.00cm \setlength{\topmargin}{0.4cm}

%%%%%%%%%%%%%%%%%%%%%%%%%%%%%%%%%%%%%%%%%%%%%%%%%%%%%%%%%%%%%%

\newtheorem{lem1}{Lema}%[chapter]
\newtheorem{th1}{Teorema}%[chapter]
\newtheorem{prop1}{Proposici\'{o}n}%[chapter]
\newtheorem{defi1}{Definici\'{o}n}%[chapter]
\newtheorem{coro1}{Corolario}%[chapter]
\newtheorem{ob1}{Observaci\'{o}n}%[chapter]
\newtheorem{obs1}{Observaci\'{o}n}%[chapter]
\newtheorem{nota1}{Nota}%[chapter]
\newtheorem{eje1}{Ejemplo}%[chapter]
\newtheorem{ejes1}{Ejemplo}%[chapter]
\newtheorem{ejerc1}{Ejercicio}%[chapter]
\newtheorem{ejercs1}{Ejercicio}%[chapter]
\newenvironment{prop}{\begin{prop1}\rm}{\end{prop1}}
\newenvironment{defi}{\begin{defi1}}{\end{defi1}}
\newenvironment{coro}{\begin{coro1}}{\end{coro1}}
\newenvironment{th2}{\begin{th1}}{\end{th1}}
\newenvironment{nota}{\begin{nota1}\rm}{\end{nota1}}
\newenvironment{lem}{\begin{lem1}}{\end{lem1}}
\newenvironment{ob}{\begin{ob1}\rm}{\end{ob1}}
\newenvironment{obs}{\begin{obs1}\rm}{\end{obs1}}
\newenvironment{eje}{\begin{eje1}\rm}{\end{eje1}}
\newenvironment{ejes}{\begin{ejes1}\rm}{\end{ejes1}}
\newenvironment{ejerc}{\begin{ejerc1}\rm}{\end{ejerc1}}
\newenvironment{ejercs}{\begin{ejercs1}\rm}{\end{ejercs1}}
%%%%%%%%%%%%%%%%%%%%%%%%%%%%%%%%%%%%%%%%%%%%%%%%%%%%%%
%\DeclareSymbolFont{cucu}{U}{bbmss}{m}{n}
\DeclareMathOperator{\CC}{\mathbb C}%{\mathalpha}{cucu}{67}
\DeclareMathOperator{\NN}{\mathbb N}%{\mathalpha}{cucu}{78}
\DeclareMathOperator{\RR}{\mathbb R}%{\mathalpha}{cucu}{82}
\DeclareMathOperator{\QQ}{\mathbb Q}%{\mathalpha}{cucu}{81}
\DeclareMathOperator{\ZZ}{\mathbb Z}%{\mathalpha}{cucu}{90}
%%%%%%%%%%%%%%%%%%%%%%%%%%%%%%%%%%%%%%%%%%%%%%%%%%%%%%
\DeclareMathOperator{\sen}{sen}
\DeclareMathOperator{\arcsen}{arcsen}
\DeclareMathOperator{\senh}{senh} \DeclareMathOperator{\fr}{Frac}
%%%%%%%%%%%%%%%%%%%%%%%%%%%%%%%%%%%%%%%%%%%%%%%%%%%%
\newcommand{\ran}[1]{{\em Im\/}(#1)}
\newcommand{\ca}[1]{{\em card\/}(#1)}
\newcommand{\sop}[1]{{\em supp\/}(#1)}
\newcommand{\tra}[1]{{\em traza\/}(#1)}
\newcommand{\dis}[1]{d (#1)}
\newcommand{\inte}[1]{{\em Int\/}(#1)}
\newcommand{\conv}[1]{{\em conv\/}(#1)}
\newcommand{\too}{\longrightarrow}
\newcommand{\oo}{\overline}
\newcommand{\cq}{\begin{flushright} $\Box$ \end{flushright}}
\newcommand{\cqd}{\, \, \rule{7pt}{7pt}}
\def\limsup{\mathop{\overline{\rm lim}}}
\def\liminf{\mathop{\underline{\rm lim}}}
\renewcommand{\theenumii}{\arabic{enumii}}
\renewcommand{\labelenumii}{\theenumi.\theenumii.}
\renewcommand{\theenumiii}{\arabic{enumiii}}
\renewcommand{\labelenumiii}{\theenumi.\theenumii.\theenumiii}
%%%%%%%%%%%%%%%%%%%%%%%%%%%%%%%%%%%%%%%%%%%%%%%%%%%%%%%%%%%%%%%%%
\begin{document}

\noindent UNA COTA PARA LOS CEROS DE LOS POLINOMIOS DE SOBOLEV

\noindent Producto Sobolev continuo: caso secuencialmente dominado.

\noindent Estimaci\'{o}n de la cota de acotaci\'{o}n de ceros de polinomios ortogonales de Sobolev.

\noindent Consideremos el siguiente producto de Sobolev:
$$
  <p(t),q(t)>=  \int_{a}^{b} p(t)q(t)w_0(t)\; dt +  \int_{a}^{b} p'(t)q'(t)w_1(t)\; dt
$$

\noindent donde $w_0(t),w_1(t)$ son pesos continuos de modo que est\'{a}n sequencialmente dominados en el sentido de que: Existe una constante $C>0$ tal que:
$$
  w_1(t) \leq Cw_0(t) , \qquad t\in [a,b]
$$

\noindent Puesto que $\dfrac{w_1(t)}{w_0(t)}$ es una funci\'{o}n continua en $[a,b]$ consideremos
$$
  \alpha= \max \{ \dfrac{w_1(t)}{w_0(t)}, t\in [a,b]\}.
$$

\noindent Consideremos el operador multiplicaci\'{o}n por $t$
$$
  \mathcal{D}: \mathbb{P}[t] \to \mathbb{P}[t], \quad \mathcal{D}(p(t))=tp(t).
$$

\noindent Veamos que $\mathcal{D}$ es acotado en este caso adem\'{a}s:
$$
  \Vert \mathcal{D} \Vert \leq C= (R^{2}+2\alpha)^{1/2}.
$$

$$
  \Vert \mathcal{D}(p(t))\Vert^{2} = \int_{a}^{b} t^2p^2(t)w_0(t)\; dt +  \int_{a}^{b} (tp'(t))^2w_1(t)\; dt= \int_{a}^{b} t^2p^2(t)w_0(t)\; dt +  \int_{a}^{b} (p(t)+tp'(t))^2w_1(t)\; dt \leq
$$
\noindent utilizando que $(a+b)^2\leq 2(a^2+b^2)$ si $a,b\geq 0$ se tiene que:
$$
  \int_{a}^{b} t^2p^2(t)w_0(t)\; dt+2\int_{a}^{b} p(t)^2w_1(t)\; dt +\int_{a}^{b} +t^2p'(t)^2w_1(t)\; dt \leq
$$
$$
  R^{2}\int_{a}^{b} p^2(t)w_0(t)\; dt +  \int_{a}^{b} (p'(t)^2w_1(t)\; dt)+ 2\int_{a}^{b} p(t)^2w_1(t)\; dt \leq
$$
$$
  R^{2} \Vert p(t)\Vert^{2} + 2\alpha \int_{a}^{b} p(t)^2w_0(t)\; dt \leq (R^{2}+2\alpha)\Vert p(t) \Vert^{2}
$$

\noindent Por ejemplo, si $[a,b]=[-1,1]$ y $\alpha =1$ quedar\'{\i}a $C=\sqrt{1+2}=\sqrt{3}$.

\begin{enumerate}
  \item Experimentar con ejemplos sequencialmente dominados $w_0(t)=t^n$ y $w_1(t)=t^{2n}$ con $n$ par en el intervalo $[-1,1]$, se aprecia que la c ota es buena ???
  \item Experimentar con ejemplos que no sean sequencialmente dominados $w_0(t)=t^{2n}$ y $w_1(t)=t^{n}$ con $n$ par en el intervalo $[-1,1]$, se aprecian resultados an\'{a}logos al caso anterior.
  \item Observar que en los casos anteriores todos los ceros están en el círculo cerrado de centro el origen y radio $\sqrt{3}$.
\end{enumerate}
\noindent
\end{document}
